\subsection{An ordinary reflection target need not carry a five-primary canonical detector}

\paragraph{Finding A36 (an actual failure of the forced-five shortcut).}
There is a constant, pointed cyclic etale cover of degree five of the
rank-one curve in A27 whose reflection quotient $E/\mathbf F_{256}$
is an ordinary genus-two curve with the Rosati-invariant
$\mathcal O_5$ action of A33 and
\[
 \#J_E(\mathbf F_{256})=288^2=2^{10}3^4.
 \tag{NF.1}
\]
Consequently every class $\kappa_R=[K_E-2R]$ with
$R\in E(\mathbf F_{256})$ has zero five-primary component.
Each such fixed class continues to have zero five-primary component
after every constant extension. The same is true of the two
correspondence constants described below. This disproves an inference
that the quintic transfer or its integral real multiplication must
force a five-primary canonical detector. It does not exclude a
three-primary detector, another useful odd prime on another target,
or five-primary classes belonging to points that first appear over
larger fields.

\begin{proof}
\emph{1. The finite quotient and its character polynomial.}
Use the field and reciprocal model already used in A32:
\[
 \begin{split}
 k&=\mathbf F_2[z]/(z^8+z^4+z^3+z+1),\qquad u=188,\quad u^2+u+1=0,\\
 C/k&:\quad Y^2+X^2Y=uX^5+X^3+uX,\qquad W=\infty.
 \end{split}
 \tag{NF.2}
\]
The integer encoding records binary coefficients in increasing powers
of $z$. The two elliptic Frobenius factors from A27 have fourth power
sums $17,-16$. Thus
\[
 \begin{aligned}
 \#C(k)=256,
 \qquad \#J_C(k)&=(257-17)(257+16)=65520,\\
                       65520&=5\cdot13104,
                       \qquad 5\nmid13104.
 \end{aligned}
 \tag{NF.3}
\]
The reduced Mumford class
\[
 \Delta=((224,80,1),(180,59))
 \tag{NF.4}
\]
has exact order five; coefficient lists in (NF.4) are in increasing
degree. Define a surjective homomorphism
$b:J_C(k)\to\mathbf Z/5$ by
$[13104]A=b(A)\Delta$. In residue order $0,1,2,3,4$, the
occupancies of $R\mapsto b([R-W])$ are
\[
                         (76,45,45,45,45).
 \tag{NF.5}
\]
The retained finite certificates specify the complete five-element
subgroup and every one of the $256$ point images. The extended-GCD
Cantor calculation independently reproduces the binary-linear
interpolation calculation, using no interpolation arithmetic.
Both preserve the curve equation during divisor reduction; the
invariant and reduction formulas are those explained in A32.

The Lang quotient for $b$, pulled back along the Abel map based at
$W$, gives a geometrically connected constant cyclic cover
$\pi:D\to C$ of degree five with a rational lift $P$ of $W$.
Indeed the Jacobian isogeny has connected domain and kernel
$\mathbf Z/5$, and the Abel map induces the required surjection
on the geometric abelian prime-to-two fundamental group, as in A12.
Its fibre above the identity, hence the fibre above $W$, is split.
For a faithful character $\chi(1)=\zeta_5$, the genus-two
Riemann--Roch calculation in recovered A1 B.2 gives
\cite[B.2]{A1run}
\[
 L_\chi(T)=1+c_\chi T+256T^2,
 \qquad
 c_\chi=76+45\sum_{i=1}^4\zeta_5^i=31.
 \tag{NF.6}
\]
Here the top coefficient is $256$ because $W$ is Weierstrass, so
$K_C\sim2W$. The same formula holds for all four faithful
characters. For the whole new factor
$\mathcal P=\ker(\operatorname{Nm}_\pi)^0$, character
factorization therefore gives
\[
                H_{\mathcal P}(T)=(1+31T+256T^2)^4.
 \tag{NF.7}
\]
Its reduction modulo two has degree four, the dimension of
$\mathcal P$, so $\mathcal P$ is ordinary. This uses the
standard Frobenius-polynomial and geometric rank comparison
already specified in A32
\cite[Propositions~2.2.4--2.2.5]{KatzCongruence}, not a
rational-class-number inference.

\emph{2. A reflection defined over the stated field.}
The hyperelliptic involution of $C$ lifts to reflections on $D$,
as in A33. Among these lifts there is a unique reflection $\tau$
fixing the chosen $k$-rational point $P$ above $W$: on the split
five-point fibre, the reflections act by $j\mapsto a-j$.
Galois preserves the cover and $P$, so this unique lift is
$k$-defined. Its quotient $v:D\to E=D/\langle\tau\rangle$
and the deck generator $\sigma$ are consequently defined over $k$.
The complete genus and isogeny calculation of A33 gives
\[
 \begin{aligned}
 g(E)&=2,\qquad f_E=2,\qquad \mathcal P\sim_k J_E\times J_E,\\
 T&=v_*\sigma v^*,\qquad T^2+T-1=0\quad\hbox{on }J_E.
 \end{aligned}
 \tag{NF.8}
\]
The embedding generated by $T$ is integral and fixed by Rosati.
Because the isogeny in (NF.8) is defined over $k$,
$H_E(T)^2=H_{\mathcal P}(T)$. Both sides have constant
coefficient one, so in the integral polynomial ring (NF.7) implies
\[
 H_E(T)=(1+31T+256T^2)^2
      =1+62T+1473T^2+15872T^3+65536T^4.
 \tag{NF.9}
\]
Evaluation at one proves (NF.1). Every $k$-rational canonical
class is an element of this finite group. Its order is therefore
prime to five, and the order of a fixed geometric point of an
abelian variety does not change upon extending its field of
definition. This proves the asserted persistent scope.
\end{proof}

\paragraph{The exact inhomogeneous correspondence identity.}
The degree-zero relation in (NF.8) is not the whole relation on
the Picard group. Work on $X=E^{(1)}$, carrying the correspondence
and the quotient map over by Frobenius twist; all the data remain
defined over $k$. Write $B$ for the locally exact differential line
bundle on $X$ attached to $E\to X$, the actual Raynaud theta
characteristic of degree one. Let $F\in\operatorname{Pic}^5(X)$
be the fibre class of the degree-five map
$X\to C^{(1)}/\langle\iota\rangle\simeq\mathbf P^1$.
Use $T$ also for the degree-doubling correspondence on the full
Picard group. For every $L\in\operatorname{Pic}^d(X)$ one has
\[
                       (T^2+T-1)L=dF.
 \tag{NF.10}
\]
Here the right side denotes addition of divisor classes, not a
degree-zero endomorphism.

For completeness, put $T_2=v_*\sigma^2v^*$. The exact
pullback--norm identity $v^*v_*=1+\tau$ and
$\sigma\tau\sigma=\tau$ give
$T^2=T_2+2$ on the Picard group. At a point $R=v(Q)$,
the divisor $R+T(R)+T_2(R)$ is the degree-five fibre over the
image of $\pi(Q)$ in the hyperelliptic projective line. On an
unramified fibre its five terms are $v(\sigma^jQ)$ for
$j=0,\ldots,4$. At a branch fibre the repeated images have
exactly their multiplicities: one point has multiplicity one and
the other two have multiplicity two. Thus the identity holds
without dividing by two. Extending additively to divisors and
then to line-bundle classes proves (NF.10).

Put
\[
             a=T(B)-2B\in J_X(k),
             \qquad b_0=F-5B\in J_X(k).
 \tag{NF.11}
\]
Substituting $L=B$ in (NF.10), and using $T(B)=2B+a$,
gives the exact identities
\[
       b_0=(T+3)a,
       \qquad (2-T)b_0=5a,
       \qquad (2-T)(T+3)=5\quad\hbox{on }J_X.
 \tag{NF.12}
\]
In particular the two constants have the same support away from
five: both $2-T$ and $T+3$ are invertible on every finite
$\ell$-primary group for $\ell\ne5$, by the last identity.
These linear relations alone do not produce a nonzero five-primary
part. In the actual example above the Frobenius twist has the same
number of $k$-rational Jacobian points as $E$, and (NF.1) shows
that the five-primary parts of both $a$ and $b_0$ are zero.
They remain zero for these fixed constants after all extensions.

Thus the proposed inhomogeneous constants do not repair the
forced-five inference. The good-prime criterion A29 still has a
separate possible three-primary application on this target, and
no absence of such a detector is asserted here.
